\documentclass{article}
\usepackage{import}
\usepackage{tcolorbox}
\usepackage{amsmath}
\usepackage{amssymb}
\usepackage{amsthm}
\usepackage{geometry}
\usepackage{enumitem}
\usepackage{pdfpages}
\usepackage{transparent}
\usepackage[hidelinks]{hyperref}
\usepackage{booktabs}
\usepackage{bm}
\usepackage{tabularx}
\tcbuselibrary{theorems}
\tcbuselibrary{skins}

% tabularx Settings
\renewcommand\tabularxcolumn[1]{>{\centering\arraybackslash}m{#1}}

% Page Settings
\newgeometry{vmargin={15mm}, hmargin={15mm, 15mm}}
\setlength{\parindent}{0pt}
\setlength{\parskip}{1em plus 1pt minus 1pt}

% hyperref Settings
\hypersetup{
    colorlinks=true,
    linkcolor=blue,
    filecolor=magenta,      
    urlcolor=cyan,
}

% inkscape Settings
\newcommand{\incfig}[1]{%
    \def\svgwidth{\columnwidth}
    \import{./figure/}{#1.pdf_tex}
}

% Definition Box Styles
\newtcbtheorem[auto counter, number within=subsection]{defn}{Definition}{
% basic settings
    theorem name and number,
    separator sign={},
    description delimiters={(}{)},
    enhanced jigsaw,
    sharp corners,
% box margins and rules
    boxrule=1pt,
    left=0.2cm,right=0.2cm,top=0.2cm,
    toptitle=0.1cm,
    bottomtitle=0.08cm,
% box colors
    colframe=white!25!black,
    colback=white,
    coltitle=white,
% box fonts
    fonttitle=\bfseries,fontupper=\normalsize
}{defn}

% Theorem Box Styles
\newtcbtheorem[auto counter, number within=subsection]{thm}{Theorem}{
% basic settings
    theorem name and number,
    separator sign={},
    description delimiters={(}{)},
    enhanced jigsaw,
    sharp corners,
% box margins and rules
    boxrule=1pt,
    left=0.2cm,right=0.2cm,top=0.2cm,
    toptitle=0.1cm,
    bottomtitle=0.08cm,
% box colors
    colframe=white!25!black,
    colback=white,
    coltitle=white,
% box fonts
    fonttitle=\bfseries,fontupper=\normalsize
}{thm}

% Lemma Box Styles
\newtcbtheorem[number within=tcb@cnt@thm]{lem}{Lemma}{
% basic settings
    theorem name and number,
    separator sign={},
    description delimiters={(}{)},
    enhanced jigsaw,
    sharp corners,
% box margins and rules
    boxrule=1pt,
    left=0.2cm,right=0.2cm,top=0.2cm,
    toptitle=0.1cm,
    bottomtitle=0.08cm,
% box colors
    colframe=white!50!black,
    colback=white,
    coltitle=white,
% box fonts
    fonttitle=\bfseries,fontupper=\normalsize
}{lem}

% Corollary Box Styles
\newtcbtheorem[number within=tcb@cnt@thm]{cor}{Corollary}{
% basic settings    
    theorem name and number,
    separator sign={},
    description delimiters={(}{)},
    enhanced jigsaw,
    sharp corners,
% box margins and rules
    boxrule=1pt,
    left=0.2cm,right=0.2cm,top=0.2cm,
    toptitle=0.1cm,
    bottomtitle=0.08cm,
% box colors
    colframe=white!50!black,
    colback=white,
    coltitle=white,
% box fonts
    fonttitle=\bfseries,fontupper=\normalsize
}{cor}

% Remark Box Styles
\newtcbtheorem[number within=subsection]{rem}{Remark}{
% basic settings
    theorem name and number,
    separator sign={},
    description delimiters={(}{)},
    enhanced jigsaw,
    sharp corners,
% box colors
    colframe=white!60!black,
    colback=white,
    coltitle=black,
% box margins and rules
    boxrule=1pt,
    left=0.2cm,right=0.2cm,top=0.2cm,
    toptitle=0.1cm,
    bottomtitle=0.08cm,
% box fonts
    fonttitle=\bfseries,fontupper=\normalsize
}{rem}

% Proof Box Styles

\title{Preliminaries of Machine Learning}
\author{Hyoung Joon, Kim}
\date{\today}

\begin{document}

\maketitle
\tableofcontents
\newpage

\section{Introduction}

This is some notes about preliminaries of machine learning.

\section{Probability}
\subsection{Basic Definitions}

% definition of probability space
\begin{defn}{Probability Space}{PSpace}
    A \emph{probability space} is a triple \(( \Omega, \mathcal{F}, \mathbb{P})\) where
    \begin{itemize}
        \item \(\Omega\) is a non-empty set called the \emph{sample space},
        \item \(\mathcal{F}\) is a \(\sigma\)-algebra of subsets of the sample space \(\Omega\), and
        \item \(\mathbb{P} :\mathcal{F} \to [0, 1]\) is a probability measure.
    \end{itemize}
\end{defn}

\subsection{Lists of Probability Distributions}

\subsubsection{Discrete Distributions}

\begin{center}
%\begin{tabular}
\begin{tabularx}{\textwidth}{c X c c c}
    \toprule
    \textbf{Name} & \textbf{Support} & \textbf{Parameters} & \textbf{PMF} & \textbf{Mean, Variance} \\
    \midrule
    Bernoulli & \(\{0, 1\}\) & \(p \in \{0, 1\}\) & Ber(\(x|p\)) = \(p^x (1-p)^{1-x}\) & \(\mathbb{E}[X] = p, \mathbb{V}[X] = p(1-p)\) \\
    \midrule
    Binomial & \(\{0, 1, 2, \ldots, N\}\) & \(N, p \in [0,1]\) & Bin(\(x|N,p\)) = \(\binom{N}{x} p^x (1-p)^{N-x}\) & \(\mathbb{E}[X] = Np, \mathbb{V}[X] = Np(1-p)\) \\
    \midrule
    Categorical & \(\{0, 1, 2, \ldots, N\}\) & \(\boldsymbol{\theta}\) & Cat(\(x|\boldsymbol{\theta}\)) = \(\prod_{k=1}^{K}\theta_{k}^{\mathbb{I}(x=k)}\) \\
    \midrule
    Multinomial & \(\bm{x}\) where \(\sum_{k = 1}^{K} x_k = N\) & \(N, \boldsymbol{\theta}\) & Mult(\(\bm{x}|\boldsymbol{\theta}\)) = \(\binom{N}{x_1, x_2, \ldots, x_K}\prod_{k = 1}^{K}\theta_{k}^{x_k}\) \\
    \midrule
    Poisson & \(\{0, 1, 2, \ldots\}\) & \(\lambda > 0\) & Pois(\(x|\lambda\)) = \(e^{-\lambda}\dfrac{\lambda^x}{x!}\) & \(\mathbb{E}[X] = \lambda, \mathbb{V}[X] = \lambda\) \\
    \bottomrule
\end{tabularx}
%\end{tabular}
\end{center}


\subsubsection{Continuous Distributions}

\begin{center}
\begin{tabularx}{\textwidth}{c c c c X}
    \toprule
    \textbf{Name} & \textbf{Support} & \textbf{Parameters} & \textbf{PDF} & \textbf{Mean, Variance} \\
    \midrule
    Gaussian &
    \(\mathbb R\) &
    \(\mu, \sigma\) &
    \(\mathcal N(\mu, \sigma^2)=\dfrac{1}{\sqrt{2\pi\sigma}}e^{-\frac{1}{2\sigma^2}(x-\mu)^2}\) &
    \(\mathbb{E}[X] = \mu, \mathbb{V}[X] = \sigma^2\)
    \\
    \midrule
    Student t-distribution &
    \(\mathbb R\) &
    \(\mu, \sigma\) &
    \(\mathcal T_\nu(x|\mu, \sigma^2)=\dfrac{\Gamma(\frac{\nu+1}{2})}{\sqrt{\pi\nu}\Gamma(\frac{\nu}{2})}\left(1+\dfrac{x^2}{\nu}\right)^{-\frac{\nu+1}{2}}\) &
    \(\mathbb E[X] = 0\) for \(\nu>0\), \(\mathbb V[X] = \frac{\nu}{\nu - 2}\) for \(\nu > 2 \)
    \\
    \midrule
    Cauchy &
    \(\mathbb R\) &
    \(\mu,\lambda\) &
    \(\mathcal C(x|\mu,\lambda)=\dfrac{1}{\lambda\pi(\frac{x-\mu}{\lambda})^2}\) &
    undefined
    \\
    \midrule
    Laplace &
    \(\mathbb R\) &
    \(\mu, b\) &
    \(\text{Laplace}(x|\mu, b) = \dfrac{1}{2b}\exp\left(-\dfrac{|x-\mu|}{b}\right)\) &
    \(\mathbb E[X] = \mu, \mathbb V[X] = 2b^2\)
    \\
    \bottomrule
\end{tabularx}
\end{center}

\end{document}